% ===========================================================================
\chapter{\label{chap:intro}Introdução}

\sigla{API}{Application Programming Interface}
\sigla{DSR}{Design Science Research}
\sigla{EAR}{Eye Aspect Ratio}
\sigla{FPS}{Quadros por segundo}
\sigla{IC}{Intervalo de confiança}
\sigla{ITS}{Intelligent Tutoring System}
\sigla{LGPD}{Lei Geral de Proteção de Dados Pessoais}
\sigla{LLM}{Large Language Model}
\sigla{LMS}{Learning Management System}
\sigla{ONNX}{Open Neural Network Exchange}
\sigla{QP}{Questão de pesquisa}
\sigla{RAG}{Retrieval-Augmented Generation}
\sigla{SPA}{Single Page Application}
\sigla{WASM}{WebAssembly}

Parte crescente do estudo universitário acontece fora da sala de aula, diante de uma tela e sem supervisão: o estudante abre o PDF da disciplina, lê, e decide sozinho quando entendeu o bastante. Esse formato exige autorregulação, isto é, a capacidade de planejar o estudo, monitorar o próprio progresso e corrigir o rumo quando algo falha. O problema é que duas das informações de que o estudante precisa para se autorregular são justamente as mais difíceis de obter sozinho: se ele está de fato prestando atenção e se de fato compreendeu o que leu.

A primeira informação falha porque a dispersão é, em grande parte, invisível para quem se dispersa. A mente pode abandonar a tarefa enquanto os olhos continuam percorrendo o texto, e muitas vezes a pessoa só percebe o desvio depois que ele aconteceu. Some-se a isso o \estr{smartphone} ao alcance da mão, que transforma o estudo em uma alternância contínua entre tarefas, com custo para o desempenho. A segunda informação falha porque a leitura passiva produz uma sensação de fluência que não equivale a compreensão. Sem uma verificação externa, o estudante descobre as lacunas tarde, em geral na prova.

A pesquisa educacional oferece uma resposta conhecida para os dois problemas: perguntas. Responder a questões sobre o conteúdo não apenas mede a aprendizagem, mas a fortalece, e perguntas intercaladas ao longo de uma aula em vídeo chegaram a reduzir a divagação mental e a melhorar o aprendizado em comparação com a mesma aula sem perguntas \cite{SKS13}. O efeito é maior quando cada resposta recebe um retorno que explica o porquê, princípio central da avaliação formativa. O obstáculo sempre foi a produção: escrever boas perguntas para cada material é trabalhoso, e no estudo autônomo não há professor para escrevê-las. O Capítulo~\ref{chap:fund} retoma essa literatura em detalhe.

\section{O problema: duas promessas tecnológicas pouco verificadas}

Dois avanços recentes parecem remover esse obstáculo. O primeiro são os modelos de linguagem de grande porte (\estr{Large Language Models}, LLMs), capazes de gerar perguntas de múltipla escolha sobre praticamente qualquer texto. O segundo é a visão computacional executada no próprio navegador, que permite estimar, a partir da câmera de um \estr{notebook} comum, se o rosto está voltado para a tela e se há um celular no campo de visão. Juntos, eles sugerem uma ferramenta que devolve ao estudante as duas informações que lhe faltam.

Cada avanço, porém, traz uma promessa que costuma ser aceita sem medição. Os LLMs produzem texto fluente que pode ser falso ou não sustentado pela fonte, fenômeno chamado de alucinação. A resposta usual é a geração aumentada por recuperação (\estr{Retrieval-Augmented Generation}, RAG), em que o modelo gera a partir de trechos recuperados de uma base de documentos. A recuperação reduz a alucinação, mas não a elimina, o que torna necessário medir, e não presumir, com que frequência a resposta correta de uma questão gerada está de fato sustentada pelo trecho que a originou. A revisão sistemática de Kurdi et al.\ \cite{Kur20} sobre geração automática de questões para fins educacionais chega a uma conclusão nessa direção: o campo precisa de relatos experimentais mais completos e de métricas de avaliação mais padronizadas para julgar a qualidade das questões geradas.

A promessa da visão computacional é outra: a de que o monitoramento por câmera pode respeitar a privacidade. Imagens do rosto são dados pessoais e, quando processadas para extrair características faciais, podem ser enquadradas como dados biométricos, que a Lei Geral de Proteção de Dados Pessoais (LGPD) classifica como sensíveis. Ferramentas de monitoramento costumam tratar a privacidade como política, com termos de uso e promessas de descarte. O princípio de \estr{privacy by design} propõe o contrário: que a proteção esteja embutida na própria arquitetura. Executar toda a inferência no navegador realiza esse princípio, porque a imagem nunca sai do dispositivo, mas cria um problema técnico pouco discutido: o desempenho passa a depender do \estr{hardware} de cada estudante, que é heterogêneo e, muitas vezes, modesto.

\section{Lacuna e proposta}

A lacuna que este trabalho endereça, portanto, não é a ausência de geradores de questões nem de detectores de engajamento, ambos amplamente estudados. É a ausência de evidência sobre o funcionamento conjunto dessas técnicas nas condições reais do estudo autônomo: material trazido pelo próprio estudante, que ninguém revisa; computadores comuns, sem placa de vídeo dedicada; e a exigência de que nenhuma imagem trafegue pela rede. Nessas condições, a qualidade das questões não pode depender de revisão humana, o desempenho não pode pressupor \estr{hardware} potente, e a privacidade precisa ser uma propriedade verificável do sistema, não uma declaração.

Para investigar essa lacuna, este trabalho constrói e avalia o \ferramenta, uma ferramenta de apoio ao estudo autônomo. O nome resume a ideia central do artefato: ancorar as perguntas no material do estudante e ancorar a atenção do estudante no material. No \ferramenta, o próprio estudante envia os PDFs de uma disciplina; o sistema fragmenta o texto, indexa os fragmentos em um banco vetorial e gera questões de múltipla escolha, cada uma a partir de um fragmento específico, que fica guardado junto à questão como sua fonte. Durante a sessão de estudo, a câmera é processada localmente para estimar um indicador de engajamento visual, o \fs, e para contar eventos de uso do celular. Um painel devolve ao estudante o cruzamento entre engajamento e desempenho, como instrumento de automonitoramento.

É importante delimitar desde já o que o \ferramenta não é. Ele não é um Ambiente Virtual de Aprendizagem no sentido institucional do termo, como Moodle ou Canvas, que integram gestão de turmas, papéis de professor e aluno, entrega de atividades e comunicação. Não há professor no sistema. O \ferramenta ocupa o espaço que esses ambientes deixam descoberto: o momento em que o estudante, sozinho, estuda o material.

\section{Questões de pesquisa}

A pergunta que orienta o trabalho pode ser enunciada em uma frase:

\begin{quote}
\emph{É viável, em computadores comuns de estudantes, combinar a geração de questões formativas ancoradas no material de estudo com a medição de engajamento visual executada inteiramente no navegador, e com que qualidade, desempenho e garantia de privacidade essa combinação funciona?}
\end{quote}

Essa pergunta se desdobra em quatro questões de pesquisa (QPs), cada uma ligada a uma das promessas discutidas acima:

\begin{itemize}
  \item \textbf{QP1 (qualidade das questões):} em que proporção as questões geradas a partir de fragmentos do material são ancoradas no fragmento, corretas e bem construídas, e quanto dessa qualidade se deve à ancoragem, em comparação com gerar sem o fragmento-fonte?
  \item \textbf{QP2 (viabilidade no navegador):} com que desempenho a detecção facial e a detecção de celular executam no navegador em \estr{hardware} heterogêneo, e com que frequência o sistema precisa recorrer ao modo degradado?
  \item \textbf{QP3 (validade do indicador):} o \fs e o contador de eventos de celular respondem como esperado a estados comportamentais conhecidos, e quão sensíveis são aos parâmetros escolhidos?
  \item \textbf{QP4 (privacidade verificável):} é possível demonstrar, por um procedimento automatizado e repetível, que nenhuma imagem da câmera deixa o navegador durante uma sessão de estudo?
\end{itemize}

As QPs foram escolhidas antes do método, a partir dos riscos que a proposta assume, e não a partir do que é mais fácil de medir. Por isso, nenhuma delas pergunta se a ferramenta melhora o aprendizado. Essa pergunta é importante, mas exige um estudo com participantes e medidas de retenção que está fora do escopo deste trabalho; o Apêndice~\ref{ap:usuarios} descreve como ele pode ser conduzido a partir dos resultados obtidos aqui.

\section{Contribuições}

Este é um trabalho de proposta de artefato com avaliação técnica, conduzido segundo a \estr{Design Science Research} (DSR) \cite{HMPR04}. Suas contribuições são:

\begin{enumerate}
  \item \textbf{Uma arquitetura de apoio ao estudo autônomo} que integra geração de questões a partir de fragmentos do material, estimativa de engajamento visual no cliente e um painel de automonitoramento, com degradação graciosa como decisão de projeto: quando o computador não suporta o detector de celular, a sessão continua apenas com os indicadores faciais.
  \item \textbf{Um protocolo de avaliação técnica reproduzível e sem participantes humanos}, baseado em um corpus fixo de materiais, em um vídeo roteirizado e rotulado que substitui a câmera, e em uma página de \estr{benchmark} que pode ser executada em qualquer navegador.
  \item \textbf{Um conjunto de métricas de qualidade para geradores de questões} que complementa a inspeção humana com verificações automáticas: ancoragem lexical, vieses de posição e de extensão da alternativa correta, duplicatas, falhas de formato e custo por questão.
  \item \textbf{A privacidade tratada como propriedade testável}, por meio de um teste automatizado que audita todo o tráfego de rede de uma sessão.
  \item \textbf{Uma análise explícita dos limites da abordagem}, em particular o fato de o \fs ser um indicador indireto de atenção e de o detector de celular depender do \estr{hardware}.
\end{enumerate}

\section{Organização do texto}

O Capítulo~\ref{chap:fund} revisa a literatura em torno de quatro controvérsias: o valor das perguntas para a atenção e a aprendizagem, a confiabilidade das questões geradas por LLMs, a validade da medição automática de engajamento e a privacidade no monitoramento por câmera. O Capítulo~\ref{chap:ferramenta} descreve o \ferramenta, com a definição formal de seus indicadores. O Capítulo~\ref{chap:metodo} apresenta o protocolo de avaliação de cada QP. O Capítulo~\ref{chap:resultados} reporta os resultados, o Capítulo~\ref{chap:discussao} os interpreta e discute as limitações, e o Capítulo~\ref{chap:conclusao} retoma a pergunta inicial e propõe trabalhos futuros.
